\chapter{Doctoral Entrance Problems}
\label{ch:problems}

The following problems are reproduced in their source language and routed by individual mathematical content.
\newpage

\section{Université Abdelhamid Ibn Badis --- Mostaganem}
\begin{problem}{(Mostaganem) — Épreuve N°01}{ca:mostaganem-2025-e1-p1}
Pour chaque $n\in\mathbb{N}^{*}$, soit $\gamma_n$ l'ellipse d'équation $4x^{2}+y^{2}-4n^{2}=0$, parcourue dans le sens direct. Calculer
$$I_n=\oint_{\gamma_n}\frac{dz}{\bigl(z^{2}-e(1+i)z+ie^{2}\bigr)^{2}}.$$
\source{\emph{\small Université Abdelhamid Ibn Badis - Mostaganem, 2025, Épreuve~n\textsuperscript{o}~01.}}
\end{problem}
\newpage
\begin{problem}{(Mostaganem) — Épreuve N°01}{ca:mostaganem-2025-e1-p2}
Soient $f$ une fonction analytique dans le disque $|z|<\rho$ et $|a|<R<\rho$. Démontrer
$$f(a)=\frac{1}{2\pi i}\oint_{|z|=R}\frac{R^{2}-a\bar a}{(z-a)(R^{2}-z\bar a)}f(z)\,dz.$$
En déduire la formule de Poisson.
\source{\emph{\small Université Abdelhamid Ibn Badis - Mostaganem, 2025, Épreuve~n\textsuperscript{o}~01.}}
\end{problem}
\newpage
\begin{problem}{(Mostaganem) — Épreuve N°01}{ca:mostaganem-2025-e1-p3}
Montrer que $f(z)=\dfrac{\cos z}{\sin \pi z}$ est méromorphe, déterminer ses pôles et résidus, puis établir son développement de Mittag-Leffler.
\source{\emph{\small Université Abdelhamid Ibn Badis - Mostaganem, 2025, Épreuve~n\textsuperscript{o}~01.}}
\end{problem}
\newpage

\section{Université Mohamed Chérif Messaadia --- Souk Ahras}
\begin{problem}{(Souk Ahras) — Épreuve N°01}{ca:souk-ahras-2025-e1-p1}
Trouver la fonction holomorphe $f$ sur $\mathbb{C}$ satisfaisant
$$zf''(z)-f'(z)-4z^3f(z)=0,$$
telle que $f(0)=1$ et $f''(0)=0$.
\source{\emph{\small Université Mohamed Chérif Messaadia - Souk Ahras, 2025, Épreuve~n\textsuperscript{o}~01.}}
\end{problem}
\newpage
\begin{problem}{(Souk Ahras) — Épreuve N°01}{ca:souk-ahras-2025-e1-p2}
Calculer $I=\int_C(1+z-3\overline z)\,dz$ pour les chemins reliant $0$ à $4$: le segment, la polygonale passant par $i$, et la demi-circonférence supérieure $|z-2|=2$.
\source{\emph{\small Université Mohamed Chérif Messaadia - Souk Ahras, 2025, Épreuve~n\textsuperscript{o}~01.}}
\end{problem}
\newpage
\begin{problem}{(Souk Ahras) — Épreuve N°01}{ca:souk-ahras-2025-e1-p3}
Soit $f=u+iv$ holomorphe sur $\Omega$. Étudier l'harmonicité de $u,v$, les conditions pour que $v+iu$ soit holomorphe, et l'implication d'une partie réelle constante.
\source{\emph{\small Université Mohamed Chérif Messaadia - Souk Ahras, 2025, Épreuve~n\textsuperscript{o}~01.}}
\end{problem}
\newpage

\section{Université Djilali Liabès --- Sidi Bel Abbès}
\begin{problem}{(Sidi Bel Abbès) — Épreuve N°01}{ca:sba-2026-e1-p3}
On donne $\cos\theta_0=2/\sqrt5$ et $\sin\theta_0=1/\sqrt5$. Calculer le module et l'argument de
$$a=3i(2+i)(4+2i)(1+i),\qquad b=\frac{(4+2i)(-1+i)}{(2-i)3i}.$$
\source{\emph{\small Université Djilali Liabès - Sidi Bel Abbès, 2026, Épreuve~n\textsuperscript{o}~01.}}
\end{problem}
